\documentclass[10pt,a4paper]{amsart}
\usepackage[margin=19mm]{geometry}
\usepackage{amsmath,amssymb,mathtools}
\usepackage[scr=rsfs,cal=euler]{mathalfa}
\usepackage{xfrac}
\usepackage[unicode,hidelinks,bookmarks=false]{hyperref}
\newcommand{\ie}{i.e.\ }
\newcommand{\cf}{cf.\ }
\newcommand{\resp}{respectively\ }
\newcommand{\Subsection}{\subsection}
\providecommand{\nobreakdash}{\nobreak\hbox{-}}
\setlength{\emergencystretch}{2em}
\multlinegap0pt
\usepackage{bm}
\usepackage{bbm}
\usepackage{dsfont}
\usepackage{xspace}
\usepackage{cancel}
\usepackage{mathtools}
\usepackage{amsthm}
\makeatletter
\@ifundefined{theorem}{\newtheorem{theorem}{Theorem}}{}
\@ifundefined{lemma}{\newtheorem{lemma}[theorem]{Lemma}}{}
\makeatother
\makeatletter
\newcommand{\R}{\mathbb{R}}
\newcommand{\T}{\mathbb{T}}                    % rooted trees
\expandafter\ifx\csname assumption\endcsname\relax
\newenvironment{assumption}[2][]{%
  \renewcommand{\thenamedassumptioninner}{#2}%
  \begin{namedassumptioninner}[#1]}{\end{namedassumptioninner}}
\fi
\makeatother
\title[Range convergence and sharp one-arm asymptotics for the crit]{Range convergence and sharp one-arm asymptotics for the critical percolation cluster in high dimensions}
\author{M. Cabezas, D. Croydon, A. Fribergh, N. Kawamoto}
\date{Source: arXiv:2607.24561v1 (CC BY 4.0)}
\begin{document}
\maketitle

\noindent\emph{Source and licence.} Range convergence and sharp one-arm asymptotics for the critical percolation cluster in high dimensions. Version arXiv:2607.24561v1 (CC BY 4.0), distributed under the Creative Commons Attribution 4.0 International licence, as recorded in the arXiv licence metadata for this exact version and shown on the abstract page https://arxiv.org/abs/2607.24561v1. Full rights notice: licenses/arxiv-2607.24561-CC-BY-4.0.txt.\par\smallskip

\noindent\textit{Editorial scope.} Counted unit: the Lemma whose source label is \texttt{lem:envelope} in the article (source0.tex lines 1275--1295, source label \texttt{lem:envelope}), delivered together with the article's own complete original proof of it (source0.tex lines 1297--1390, one \texttt{proof} environment of 94 lines). No mathematics is written, repaired or re-derived: statement and proof are the article's own text line for line. Required context retained uncounted: eq:IT (lines 1243--1249). Every definition, display and label of those lines is retained unchanged. Omitted from this package and not redistributed: 1--1242, 1250--1274, 1296--1296, 1391--3267. Nothing inside a retained excerpt is omitted or altered: every retained body is delivered exactly as the article's lines are written, so no omission receipt of any kind accompanies this package. Citations: the retained excerpts carry no citation command at all, so no bibliography entry of the article is transcribed and no reference is redistributed. Reference adaptation: none was needed and none was made; every reference inside the delivered unit resolves inside it, so no unresolved reference or citation is produced. Printed slips kept literal: no printed word, symbol, hypothesis or number of the counted statement or of its proof was altered. Layout and class: the article's own class is \texttt{article} with packages including geometry, amsmath, amssymb, amsthm, mathtools, enumitem, xcolor, hyperref, graphicx, bbm, extarrows, tikz; the wrapper is the lane's pinned \texttt{amsart} editorial wrapper at 10pt on A4, which restates the theorem-like environments the retained text needs and adds the article's own macro definitions that the retained text needs (\texttt{\textbackslash R}, \texttt{\textbackslash T}), with extra wrapper packages (bm, bbm, dsfont, xspace, cancel, mathtools). The delivered numbering is the unit's own. Assets: no retained span contains a figure environment, a graphics-inclusion command, a PDF-page inclusion, a table or a TikZ picture; no image file is copied into this package, the package declares no asset dependency and no image file is redistributed with it. Licence: version arXiv:2607.24561v1 of this article is distributed under the Creative Commons Attribution 4.0 International licence, as recorded in the arXiv licence metadata for this exact version and shown on the abstract page https://arxiv.org/abs/2607.24561v1. A full rights notice is delivered at licenses/arxiv-2607.24561-CC-BY-4.0.txt. Characters: every retained excerpt is pure ASCII, so the delivered engine, XeLaTeX with its default Latin Modern text font, reports no missing character.
\begin{equation}\label{eq:IT}
I_T(z_1,\dots,z_k)
:=\int_{(\R^d)^{V_{\mathrm{int}}(T)}}
\prod_{\{a,b\}\in E(T)}G(u_a-u_b)
\prod_{v\in V_{\mathrm{int}}(T)} du_v,
\qquad u_0:=0,\ u_i:=z_i .
\end{equation}

\begin{lemma}[Uniform domination bound]\label{lem:envelope}
Let $d>4$ and $R_0\ge1$. There is a constant $C_\star=C_\star(R_0)$$<\infty$
such that for
every $k\ge1$, every $T\in\T_k$, and all measurable
$\phi_1,\dots,\phi_k$ with $|\phi_i|\le\mathbf 1_{\overline B_{R_0}}$,
\begin{equation}\label{eq:envelope}
\int_{(\R^d)^{k}}\int_{(\R^d)^{V_{\mathrm{int}}(T)}}
\prod_{\{a,b\}\in E(T)}\bar G(u_a-u_b)
\prod_{v\in V_{\mathrm{int}}(T)}du_v\,
\prod_{i=1}^k|\phi_i(z_i)|\,dz_i
\;\le\;C_\star^{\,k},
\end{equation}
with $u_0:=0$ and $u_i:=z_i$ as in \eqref{eq:IT}. In particular
$I_T(z_1,\dots,z_k)<\infty$ for Lebesgue-a.e.\ $(z_1,\dots,z_k)$, and, after
enlarging $C_\star$ if necessary,
\[
\int_{(\R^d)^k}I_T(\vec z)\prod_i|\phi_i(z_i)|\,dz_i\le
C_\star^{\,k}.
\]

\end{lemma}

\begin{proof}
Since the integrand depends on the $\phi_i$ only through $|\phi_i|$, we may assume $0\le\phi_i\le\mathbf 1_{\overline B_{R_0}}$. By scaling we may take $R_0=1$ (rescaling space by
$R_0$ multiplies the left side of \eqref{eq:envelope} by $R_0^{2(2k-1)}$ ---
a factor $R_0^{d}$ per integration variable against $R_0^{-(d-2)}$ per edge
kernel --- which is again geometric in $k$). Set
\[
m(u):=(1+|u|)^{-(d-2)} .
\]
We will use repeatedly the elementary polar-coordinate identity
\begin{equation}\label{eq:polar}
\int_{|v|\le R}|v|^{-(d-2)}\,dv=\frac{\omega_{d-1}}{2}\,R^2
\qquad(R>0),
\end{equation}
where $\omega_{d-1}$ is the surface area of the unit sphere (the exponent
$d-1$ from the Jacobian minus $d-2$ from the kernel leaves $r^1$).
We make two claims.

\emph{Claim 1 (leaf reduction).} There is $c_1=c_1(d)$ such that for
$0\le\phi\le\mathbf 1_{\overline B_1}$,
\[
f_\phi(u):=\int\bar G(u-z)\phi(z)\,dz\;\le\;c_1\,m(u).
\]
Indeed, for $|u|\ge2$ and $z\in\overline B_1$ one has $|u-z|\ge|u|-1$, and
since $t\mapsto\frac{1+t}{t-1}$ is decreasing on $(1,\infty)$ with value $3$
at $t=2$, $(|u|-1)^{-(d-2)}\le 3^{d-2}(1+|u|)^{-(d-2)}$; thus
$\bar G(u-z)\le 3^{d-2}m(u)$, and integrating over $z\in\overline B_1$
gives $f_\phi(u)\le 3^{d-2}(\omega_{d-1}/d)\,m(u)$ (here $\omega_{d-1}/d=|B_1|$ is the volume of the unit ball). For $|u|\le2$,
$\int_{B_1}\bar G(u-z)\,dz\le\int_{B_3}|w|^{-(d-2)}dw<\infty$ (finite
by \eqref{eq:polar} with $R=3$), while $m(u)\ge(1+2)^{-(d-2)}=3^{-(d-2)}$
there, so $f_\phi(u)\le c\,m(u)$. Taking $c_1$ the larger of the two constants
proves the claim.

\emph{Claim 2 (merging).} There is $c_2=c_2(d)$ such that if
$0\le f\le\kappa\,m$ and $0\le g\le\kappa'\,m$, then
\[
h(u):=\int\bar G(u-w)\,f(w)g(w)\,dw\;\le\;c_2\,\kappa\kappa'\,m(u).
\]
To see this, note $fg\le\kappa\kappa' m^2$ and
$m^2(w)=(1+|w|)^{-2(d-2)}\in L^1(\R^d)$ because $2(d-2)>d$ for $d>4$; write
$c_m:=\|m^2\|_{L^1}$. For $|u|\le2$,
\[
h(u)\le\kappa\kappa'\Bigl[\int_{|u-w|\le1}|u-w|^{-(d-2)}\,dw
+c_m\Bigr]\le c\,\kappa\kappa'\;\le\;c'\,\kappa\kappa'\,m(u),
\]
the local integral being finite by \eqref{eq:polar} with $R=1$, and the last inequality using only that $|u|\le2$, so that $m(u)\ge3^{-(d-2)}$.
For $|u|\ge2$, split the integral into three regions:
(a) $|w|\le|u|/2$: then $\bar G(u-w)\le(|u|/2)^{-(d-2)}$, contributing
at most $2^{d-2}c_m\kappa\kappa'|u|^{-(d-2)}\le c\,\kappa\kappa'\,m(u)$ (using $|u|^{-(d-2)}\le2^{d-2}m(u)$, valid since $|u|\ge2$);
(b) $|w|\ge|u|/2$ and $|u-w|\le|u|/2$: then
$m^2(w)\le(1+|u|/2)^{-2(d-2)}$, while by the substitution $v=u-w$ and
\eqref{eq:polar},
\[
\int_{|u-w|\le|u|/2}|u-w|^{-(d-2)}\,dw
=\int_{|v|\le|u|/2}|v|^{-(d-2)}\,dv=\frac{\omega_{d-1}}{2}\Bigl(\frac{|u|}{2}\Bigr)^2
\le c|u|^2;
\]
hence this region contributes at most
$c\,\kappa\kappa'\,|u|^{2}(1+|u|/2)^{-2(d-2)}
\le c'\,\kappa\kappa'(1+|u|)^{-(d-2)}=c'\,\kappa\kappa'\,m(u)$, using $1+|u|/2\ge(1+|u|)/2$,
$|u|\le1+|u|$ and $2(d-2)-2\ge d-2$, i.e.\
$d\ge4$;
(c) $|w|\ge|u|/2$ and $|u-w|\ge|u|/2$: then
$\bar G(u-w)\le(|u|/2)^{-(d-2)}$, contributing at most
$2^{d-2}c_m\kappa\kappa'|u|^{-(d-2)}\le c\,\kappa\kappa'\,m(u)$ (again $|u|^{-(d-2)}\le2^{d-2}m(u)$ since $|u|\ge2$). As each of the three regions is bounded by a multiple of $\kappa\kappa'\,m(u)$, Claim 2 follows.

Now fix $T\in\T_k$ and orient it away from the root leaf $0$. Integrate out
the variables in order of decreasing graph distance from the root, as
follows: each non-root leaf $i$ is first replaced by the function
$f_{\phi_i}$ of its parent's position (Claim 1); then, iteratively, each
internal vertex $v$ whose two children have already been reduced to
functions $f,g$ of $u_v$ with $f\le\kappa m$, $g\le\kappa'm$ is integrated
out, producing the function $h\le c_2\kappa\kappa' m$ of its parent's
position (Claim 2). After all $k-1$ internal vertices are processed, there
remains the edge into the root. If $k=1$ the tree is the single edge
$\{0,1\}$ and there is nothing left to integrate: the left side of
\eqref{eq:envelope} equals $f_{\phi_1}(0)\le c_1\,m(0)=c_1$ by Claim~1
alone. If $k\ge2$ the root's unique neighbour is an internal vertex; once
every other internal vertex is processed, the two subtree weights
$f\le\kappa\,m$ and $g\le\kappa'\,m$ hang at this neighbour, and the final
integral
\[
\int \bar G(0-u)\,f(u)g(u)\,du
\]
is Claim 2's computation evaluated at the root, giving at most
$c_2\,\kappa\kappa'\,m(0)=c_2\,\kappa\kappa'$. Each application of
Claim 1 contributes a factor $c_1$ (there are $k$ of them) and each
application of Claim 2 a factor $c_2$ (at most $k-1$ of them), so the total
is at most $C_\star^k$ with $C_\star=C_\star(R_0)$. Tonelli's
theorem (all integrands nonnegative) justifies the iterated integration and
yields the a.e.\ finiteness of $I_T$ and \eqref{eq:envelope}. For the final
display, each of the $2k-1$ edges of $T$ carries $G=c_d\bar G$, so the left
side gains a factor $c_d^{\,2k-1}\le\bigl((1\vee c_d)^{2}\bigr)^{k}$, which is
absorbed by enlarging $C_\star$ to $(1\vee c_d)^{2}C_\star$.
\end{proof}

\end{document}
